\section*{Chapter \ref{ch:iv2:betting-and-kelly} --- Betting and Kelly}
\markboth{\protect\omTocMark{}{Answers: Betting and Kelly}}{\protect\omTocMark{}{Answers: Betting and Kelly}}

\begin{solution}{iq:iv2:betting-and-kelly:1}
At 4 to 1 you lose $1/6$ of a dollar per dollar staked, and 5 to 1 is fair. The coin bookmaker's margin is 5\% of the money it takes.

Expected profit is $\tfrac16 \times 4 - \tfrac56 \times 1 = -\tfrac16$. Fair odds make it zero, so $\tfrac16 b = \tfrac56$ and $b = 5$. In general, fair odds on an event of probability $p$ are $(1 - p)/p$ to 1.

For the coin, each decimal odd of 1.90 implies a probability of $1/1.90 \approx 52.6\%$, and the two add to about 105.3\%. If 100 goes on each side, the bookmaker takes 200 and pays $100 \times 1.90 = 190$ whichever side wins. It keeps 10 of the 200, or 5\%, with no risk.

\iqwrong{Reading 5.3\% as the margin. The excess over 100\% is measured on the payout. The share of the money taken is $1 - 1/1.0526 = 5\%$.}
\end{solution}

\begin{solution}{iq:iv2:betting-and-kelly:2}
Stake about \$418.60 on $X$ and \$581.40 on $Y$. Either way you get back about \$1,046.51, a profit of about \$46.51, or 4.65\%.

The implied probabilities are $1/2.50 = 0.40$ and $1/1.80 \approx 0.5556$, and they add to $0.9556$, below 1. That gap is the arbitrage\index{arbitrage}. Stake in proportion to the implied probabilities, so that the two payouts are equal: $0.40/0.9556 \approx 41.86\%$ on $X$ and the rest on $Y$. Each dollar then returns $1/0.9556 \approx 1.0465$.

\iqwrong{Splitting the money evenly. With \$500 on each you get \$1,250 if $X$ wins and \$900 if $Y$ wins, so you have a profitable bet on $X$ and no lock.}
\end{solution}

\begin{solution}{iq:iv2:betting-and-kelly:3}
Stake 10\% of your wealth. Log wealth then grows by about 0.0050 per bet, or about 0.5\%.

Wealth after many bets depends on the average of its logarithm, so you maximise $g(f) = p\ln(1 + f) + q\ln(1 - f)$, with $q = 1 - p$. Setting the derivative to zero gives $f = p - q = 0.10$. This is the Kelly criterion\index{Kelly criterion}, and for an even-money bet it is simply your edge. Then $g = 0.55\ln 1.1 + 0.45\ln 0.9 \approx 0.0050$. A quick estimate is $\tfrac12 f^2 = 0.005$, half the square of the edge.

\iqwrong{Maximising expected wealth, $1 + (p - q)f$. That grows with $f$, so it tells you to stake everything, and then one loss ends the game.}
\end{solution}

\begin{solution}{iq:iv2:betting-and-kelly:4}
At 40\%, stake 10\%. At 30\%, do not bet: the formula gives $-5\%$, and you are not offered the other side.

For a bet paying $b$ to 1, maximising $p\ln(1 + bf) + q\ln(1 - f)$ gives $f = p - q/b$, which is expected profit per dollar over the odds. At 40\%, the expected profit is $0.4 \times 2 - 0.6 = 0.2$, and dividing by 2 gives 0.1, the same as $0.4 - 0.6/2$. At 30\%, the expected profit is $-0.1$, so the bet is bad.

\iqwrong{``Size by the probability of winning.'' A long shot with a good price can deserve more than a likely bet at a poor price. The edge and the odds decide.}
\end{solution}

\begin{solution}{iq:iv2:betting-and-kelly:5}
Expected wealth rises 5\% a round. A typical path falls about 5.1\% a round, because the median multiplier over two rounds is $1.5 \times 0.6 = 0.9$. So if you must commit all your wealth every round, expect to lose, whatever the mean says.

The mean is $\tfrac12 \times 1.5 + \tfrac12 \times 0.6 = 1.05$. A path's growth is the average of the logarithms, $\tfrac12 \ln 0.9 \approx -0.0527$, so the typical multiplier per round is $\sqrt{0.9} \approx 0.949$. The two differ because the mean is lifted by rare paths with very many wins, on which almost nobody ends up. This gap, about half the variance of the return, is volatility drag\index{volatility drag}.

\iqwrong{``Positive expected return, so all in.'' A rise of 50\% followed by a fall of 40\% leaves you at 0.9, in either order.}
\iqfollow{What should you do with the bet?}{Stake part of your wealth. Maximising $\tfrac12\ln(1 + 0.5f) + \tfrac12\ln(1 - 0.4f)$ gives $f = 25\%$, and the growth becomes positive, $\tfrac12\ln 1.0125 \approx 0.6\%$ a round.}
\end{solution}

\begin{solution}{iq:iv2:betting-and-kelly:6}
Expected final wealth is $1.2^{10} \approx \$6.19$. The probability you finish with anything is $0.6^{10} \approx 0.6\%$. The one path in which you win all ten pays \$1,024, and it supplies the whole mean. So you would not play this way unless a single loss did not matter to you.

Each bet multiplies expected wealth by $0.6 \times 2 = 1.2$, so maximising the mean does indeed say stake all. It maximises the mean by putting the whole outcome in one path of probability 0.6\%.

\iqwrong{Treating the mean as what you can expect to have. A 99.4\% chance of nothing is what you can expect.}
\iqfollow{And at the 20\% Kelly fraction?}{Expected wealth is $1.04^{10} \approx 1.48$, you never reach zero, and even ten losses leave $0.8^{10} \approx 0.107$.}
\end{solution}

\begin{solution}{iq:iv2:betting-and-kelly:7}
The optimum is 20\%, with growth of about 0.0201 per bet. Half of it, 10\%, grows about 0.0150, which is about 75\% as fast. The standard deviation of the change in log wealth falls by about half, from about 0.199 to about 0.098.

Near the optimum, growth is close to a parabola. A multiple $c$ of the Kelly fraction keeps a share $c(2 - c)$ of the maximum growth, so half Kelly\index{fractional Kelly} keeps $\tfrac34$ for half the risk. That is the trade people accept: the first quarter of growth you give up buys half the volatility.

The standard deviation is $\sqrt{pq}\,\ln\tfrac{1 + f}{1 - f}$. It gives $0.490 \times \ln(1.2/0.8)$ for full Kelly and $0.490 \times \ln(1.1/0.9)$ for half.

\iqwrong{Thinking half the stake means half the growth. Growth is concave in the fraction, so the first half of the stake earns three quarters of it.}
\end{solution}

\begin{solution}{iq:iv2:betting-and-kelly:8}
Stake 50\% of your wealth. The growth rate is $\tfrac12\ln 1.125 \approx 0.059$ per bet.

You maximise $\tfrac12\ln(1 + f) + \tfrac12\ln(1 - \tfrac12 f)$. The derivative is zero when $1 - \tfrac12 f = \tfrac12(1 + f)$, so $f = \tfrac12$. In general, when you gain a share $b$ of the stake with probability $p$ and lose a share $a$ otherwise, $f = p/a - q/b$. Here that is $\tfrac12/\tfrac12 - \tfrac12/1 = \tfrac12$.

The usual rule $p - q/b$ assumes you can lose the whole stake, so $a = 1$. Here it would say zero, as though the bet were a fair even-money flip. The formula's stake is really the amount you can lose: staking 50\% puts 25\% of your wealth at risk.

\iqwrong{Using the expected return of 25\% as the answer. A size depends on the shape of the bet as well as on its mean.}
\end{solution}

\begin{solution}{iq:iv2:betting-and-kelly:9}
You accept with a wealth of at least \$300. The expected profit is \$25, which is positive, and it is not enough because a log investor also pays for the risk relative to wealth.

Accept if $\tfrac12\ln(w + 150) + \tfrac12\ln(w - 100) \geq \ln w$, that is, $(w + 150)(w - 100) \geq w^2$. So $50w - 15{,}000 \geq 0$, which gives $w \geq 300$.

A check by the Kelly formula: this is a bet at $b = 1.5$ with $p = \tfrac12$, so $f = \tfrac12 - \tfrac12/1.5 = \tfrac16$. Staking \$100 is Kelly at a wealth of \$600. It is twice Kelly at \$300, the point where growth reaches zero.

\iqwrong{Saying any bet with a positive mean is worth taking. At a wealth of \$110, this flip risks almost everything for a small expected gain, and log utility refuses it.}
\end{solution}

\begin{solution}{iq:iv2:betting-and-kelly:10}
The growth-optimal leverage is $\mu/\sigma^2 = 0.08/0.0256 = 3.125$. The growth above the borrowing rate is $\mu^2/(2\sigma^2) = 12.5\%$ a year. With Sharpe ratio $S = \mu/\sigma$, the leverage is $S/\sigma$, the growth is $S^2/2$, and the leveraged volatility is $S$, here 50\%.

With leverage $f$, the log of wealth grows at $g(f) = f\mu - \tfrac12 f^2\sigma^2$, the second term being volatility drag. This is a parabola with its top at $f^* = \mu/\sigma^2$. Substituting gives $g^* = \tfrac12\mu^2/\sigma^2$.

\iqwrong{Thinking the answer depends on the strategy's scale. Doubling a strategy's positions doubles $\mu$ and $\sigma$, halves the leverage, and leaves $S$, and so the growth, unchanged.}
\iqwhy{This is how Kelly appears in practice: a target volatility equal to the Sharpe ratio. A Sharpe ratio of 0.5 says at most 50\% volatility, and hardly anyone runs that much.}
\end{solution}

\begin{solution}{iq:iv2:betting-and-kelly:11}
At 40\%, growth is $0.6\ln 1.4 + 0.4\ln 0.6 \approx -0.0024$ per bet, so a typical path shrinks, despite the positive edge. Growth is zero at about 38.9\%.

The optimum is 20\%. Solving $0.6\ln(1 + f) + 0.4\ln(1 - f) = 0$ numerically gives $f \approx 0.389$.

The rule of thumb comes from the parabola $g(f) \approx (p - q)f - \tfrac12 f^2$, which is symmetric about its top and so has its second zero at $2f^*$. The true function has a logarithm that falls more sharply as $f$ approaches 1, so the right side of the curve is steeper than the left. So the real zero sits below twice the optimum, and overbetting costs more than underbetting by the same amount. At 10\%, ten points below the optimum, growth is about 0.0150. At 30\%, ten points above it, growth is about 0.0147.

\iqwrong{``More edge, more stake, more growth.'' Past the optimum you pay more in volatility drag than you earn in edge, and at about twice it you earn nothing.}
\iqwhy{An estimated edge that is twice the truth turns a strategy that should compound into one that goes nowhere, with all of the risk.}
\end{solution}

\begin{solution}{iq:iv2:betting-and-kelly:12}
Stake 25\%. The growth rate is $0.3\ln 1.125 \approx 0.0353$ per bet. The money-back outcome adds nothing to the logarithm, so it plays no part in the size, given the other two probabilities.

Maximise $0.3\ln(1 + 2f) + 0.4\ln 1 + 0.3\ln(1 - f)$. The derivative is zero when $0.6/(1 + 2f) = 0.3/(1 - f)$, so $0.6 - 0.6f = 0.3 + 0.6f$, and $f = \tfrac14$. Then $0.3\ln 1.5 + 0.3\ln 0.75 = 0.3\ln 1.125$.

Compare the same bet with no pushes, where a 50-50 win or loss pays 2 to 1. That has $f = \tfrac12 - \tfrac12/2 = 25\%$ too. The ratio of the two probabilities is what fixes the size. The growth scales with the share of the bets that settle: it is $0.6$ of that bet's $0.5\ln 1.125 \approx 0.0589$.

\iqwrong{Lowering the win probability to 30\% in the two-outcome formula, which gives $0.3 - 0.7/2 < 0$, and so no bet. The formula with a whole loss only applies when there are only two outcomes.}
\end{solution}

\begin{solution}{iq:iv2:betting-and-kelly:13}
Stake about 19.2\% on each, about 38.5\% in total. Growth per round is about 0.0395, against about 0.0403 for making the two bets one after the other at 20\%.

If both win you have $1 + 2f$, if one wins you have 1, and if both lose you have $1 - 2f$. The probabilities are 0.36, 0.48 and 0.16. The derivative of $0.36\ln(1 + 2f) + 0.16\ln(1 - 2f)$ is zero at $0.72(1 - 2f) = 0.32(1 + 2f)$, so $f = 0.4/2.08 \approx 0.1923$.

You stake less than 20\% each because the second bet cannot use the result of the first. The stakes are taken from the same wealth, and they are at risk at once. So the combined loss of 38.5\% is a bigger bite than two losses in a row at 20\%, which leave 64\%.

\iqwrong{Adding the two 20\% fractions. For small edges that is close, but it is wrong in principle, and the error grows with the number of simultaneous bets.}
\iqfollow{And with many such bets?}{The total stake rises, but below the sum of the single fractions. With 5 such bets, the optimum is about 16.5\% on each, about 83\% in total. With 50, it is about 2\% on each, just inside the limit at which losing them all would take you to zero.}
\end{solution}

\begin{solution}{iq:iv2:betting-and-kelly:14}
At correlation 0.5, each gets about 1.67 of leverage, 3.33 in total, and growth is about 16.7\%. At zero correlation, each gets 2.5, 5 in total, and growth is 25\%.

For several assets, $f^* = \Sigma^{-1}\mu$, and $g^* = \tfrac12\mu^\top\Sigma^{-1}\mu$. By symmetry the two weights are equal, so each weight solves $\sigma^2(1 + \rho)f = \mu$. That is $0.04 \times 1.5 \times f = 0.1$, so $f = 1.667$. At $\rho = 0$, $0.04f = 0.1$ gives $f = 2.5$.

The growth is half the squared Sharpe ratio of the best mix, $\tfrac12 \times 2 \times 0.1 \times f$. So $0.1667$ at $\rho = 0.5$, and $0.25$ at $\rho = 0$. The mix's Sharpe ratio is $0.5\sqrt{2/(1 + \rho)}$, about 0.577 at 0.5 and about 0.707 at 0.

\iqwrong{Adding the two single-strategy Kelly leverages of 2.5 each, which ignores the correlation. With a correlation of 0.5, that much leverage has more volatility than the mix calls for, and growth falls from about 16.7\% to 12.5\%.}
\end{solution}

\begin{solution}{iq:iv2:betting-and-kelly:15}
Stake 50\% on the tip each time. Growth is $\ln 2 - H(0.75) \approx 0.131$ per toss, where $H$ is the entropy of the tip's accuracy, in natural units. With a perfect tipster it is $\ln 2 \approx 0.693$, and with a tipster who is right half the time it is zero.

From your side the bet is even money, won with probability 75\%, so $f = 0.75 - 0.25 = 0.5$. The growth is $0.75\ln 1.5 + 0.25\ln 0.5$. Rewrite that as $\ln 2 + 0.75\ln 0.75 + 0.25\ln 0.25$, which is the formula. Kelly's 1956 paper made this link: the growth rate equals the information the tip carries about the toss.

A perfect tipster lets you stake everything and double each time. A useless one gives you no reason to bet, so growth is zero, and at 25\% accuracy you bet against the tip and have the same 0.131.

\iqwrong{Saying a 75\% tipster gives 75\% of the perfect growth. It gives about 19\% of it, because information falls quickly as accuracy leaves 100\%.}
\iqwhy{It puts a price on a signal. A signal is worth what it adds to growth, and a modest gain in accuracy near a coin flip adds little.}
\end{solution}

\begin{solution}{iq:iv2:betting-and-kelly:16}
Split your wealth 50-30-20, in proportion to the win probabilities. Growth is about 0.069 per race. The split does not change with the payouts.

With a split $w$, growth is $\sum_i p_i\ln(w_i o_i) = \sum_i p_i\ln o_i + \sum_i p_i\ln w_i$. The payouts appear only in the first term, and the $w_i$ only in the second. The second is largest at $w_i = p_i$ when the weights must add to 1, a fact called Gibbs' inequality. This is proportional betting\index{proportional betting}. At 3 for 1, growth is $\sum_i p_i\ln(3p_i) = \ln 3 - H(p) \approx 1.0986 - 1.0297 = 0.069$.

With payouts 2, 4 and 8, the split stays 50-30-20, and growth becomes $0.5\ln 1 + 0.3\ln 1.2 + 0.2\ln 1.6 \approx 0.149$.

The result depends on having to stake everything. When you can hold some back, and the bookmaker's odds imply probabilities that add to more than 1, you may keep cash and leave a bad horse out altogether.

\iqwrong{Putting everything on the favourite. That is the most likely horse, but one loss wipes you out, and growth is $-\infty$.}
\end{solution}

\begin{solution}{iq:iv2:betting-and-kelly:17}
At \$100, a \$10 bet is 10\% of your wealth, and growth is about $-0.0030$ per bet, so you should not take it. Growth reaches zero at a wealth of about \$250. \$10 is the Kelly stake at a wealth of \$500.

Kelly is $0.51 - 0.49 = 2\%$. Then $0.51\ln 1.1 + 0.49\ln 0.9 \approx -0.0030$. Growth is zero when $f \approx 0.0400$, about twice Kelly, so wealth is $\$10/0.04 \approx \$250$. And $\$10/0.02 = \$500$.

So with a small edge and a minimum stake, there is a wealth below which a good bet is a bad one. Between \$250 and \$500, a \$10 bet is oversized but still grows, and it may be worth making. Below \$250, you lose by taking it, however good its mean looks.

\iqwrong{``The edge is positive, so play.'' Expected wealth rises 20 cents a bet, but growth is negative: that is a sign you are too small to take this bet.}
\end{solution}

\begin{solution}{iq:iv2:betting-and-kelly:18}
The expected gain is \$10 per bet, and the probability of ever being wiped out is $(2/3)^{20} \approx 0.03\%$. Once your wealth passes \$250, Kelly would stake more than the limit, so your wealth grows linearly, about \$10 a bet, while Kelly's would compound.

With stakes fixed at \$50, this is gambler's ruin\index{gambler's ruin} with 20 steps of capital against an infinitely rich opponent. The chance of ever losing 20 net steps is $(q/p)^{20}$, with $q/p = \tfrac23$. The expected gain is $50 \times (0.6 - 0.4)$.

The limit changes what is optimal. With wealth well above \$250, you always stake the limit. Now your growth in log terms falls as you get richer, and risk of ruin is tiny, so you would want more room, which is a reason to look for more tables, not to change your stake.

\iqwrong{Staking 20\% of your wealth until it hits the limit and thinking the Kelly formula still holds above it. Past that point there is no decision left, only the number of bets.}
\iqfollow{Same limit, with \$200?}{At \$200, the stake of \$50 is 25\%, above Kelly but below twice it. The ruin chance is $(2/3)^4 \approx 20\%$ if you always bet the limit, so you should bet below it until you have grown.}
\end{solution}

\begin{solution}{iq:iv2:betting-and-kelly:19}
At 40\%, go all in once, for a 40\% chance. \$1 bets give about $2.5 \times 10^{-18}$. At 60\%, make small bets, for a chance within about $2.5 \times 10^{-18}$ of 1, while one all-in gives 60\%.

Each case comes from gambler's ruin: with $r = q/p$, the chance of reaching 200 from 100 in \$1 steps is $(1 - r^{100})/(1 - r^{200}) = 1/(1 + r^{100})$. At $p = 0.4$, $r = 1.5$, and $1.5^{-100} \approx 2.5 \times 10^{-18}$. At $p = 0.6$, $r = \tfrac23$, and the chance is $1 - (\tfrac23)^{100}$, near enough.

In a bad game, each bet costs you on average, so you want as few as possible: this is bold play\index{bold play}. In a good game, each bet makes you money on average, and small stakes let the law of large numbers do the work.

\iqwrong{``Kelly is optimal, so stake Kelly.'' Kelly maximises long-run growth, not the chance of reaching a target. At 40\% it says not to bet, and you would then certainly miss.}
\iqwhy{An objective changes the bet. A desk that must make a target by year end, with nothing to lose, has the same reason as this gambler to raise its risk, and this is why limits and pay structures matter.}
\end{solution}

\begin{solution}{iq:iv2:betting-and-kelly:20}
The expected loss is \$50. The largest premium you would pay at \$20,000 is about \$57.45, so at \$60 you do not buy. At \$10,000, the same formula gives about \$69.08, so you buy.

The premium $\pi$ makes you indifferent when
\[
\ln(W - \pi) = 0.99\ln W + 0.01\ln(W - 5000),
\]
so $\pi = W\left(1 - (1 - 5000/W)^{0.01}\right)$. At $W = 20{,}000$, that is $20{,}000 \times (1 - 0.75^{0.01}) \approx 57.45$. At $W = 10{,}000$, it is $10{,}000 \times (1 - 0.5^{0.01}) \approx 69.08$.

At \$20,000 you would pay up to \$7.45 above the expected loss, and this premium asks \$10 above it. The same risk is a larger share of a smaller wealth, so at \$10,000 you would pay up to \$19.08 above it. A person with much more wealth would refuse any premium much above \$50.

\iqwrong{Refusing because the expected loss is \$50, so any premium above it is a bad bet. That is true in expected value, and yet people with limited wealth should sometimes pay more than the expected loss.}
\end{solution}

\begin{solution}{iq:iv2:betting-and-kelly:21}
The mean gain is \$5,000, the standard deviation is \$1,500, and the probability of a loss is about 0.04\%. A log investor accepts one flip with a wealth of \$200 or more. The colleague is not consistent if that refusal holds at every wealth the colleague could reach.

Each flip has a mean of 50 and a standard deviation of 150. Over 100 flips, the mean is multiplied by 100 and the standard deviation by 10. You lose only with at most 33 heads, because $200h - 100(100 - h) < 0$ needs $h < 33.3$. That probability is about 0.00044. For one flip, accept if $(w + 200)(w - 100) \geq w^2$, so $w \geq 200$.

This is Samuelson's fallacy of large numbers\index{Samuelson's fallacy}. If a person with fixed preferences refuses one flip at every wealth that could occur in the sequence, then that person must refuse the whole sequence, by backward induction from the last flip. Refusing one and accepting many can be consistent only if the refusal holds at some wealths and not at others. A log investor is not an example: below \$200 they refuse one flip, and at any wealth up to \$10,000 they also refuse a commitment to 100 flips, because 100 tails would wipe them out. Above \$10,000 they accept both.

\iqwrong{``Diversification across time makes it safe.'' The stakes are added, so the gain's standard deviation grows too, by 10 times. The chance of a loss falls, but the possible size of the loss grows to \$10,000.}
\iqwhy{This argument comes up over whether long horizons justify more equities. The answer depends on the utility, and not just on the arithmetic of the average.}
\end{solution}

\begin{solution}{iq:iv2:betting-and-kelly:22}
At full Kelly, log wealth grows 0.5 a year, a typical path doubles in about 1.39 years, and volatility is 100\%. At half Kelly, it grows 0.375 a year, doubles in about 1.85 years, and has a volatility of 50\%.

Kelly leverage makes the strategy's volatility equal to $S$, here 100\%. A multiple $c$ gives volatility $cS$ and growth $(c - \tfrac12 c^2)S^2$. For $c = 1$, growth is 0.5, and for $c = \tfrac12$ it is 0.375. The doubling time is $\ln 2/g$.

So half Kelly takes a third longer to double, about 0.46 years more, for half the volatility. A desk would judge that a bargain, and that is why almost no one actually runs full Kelly.

\iqwrong{Using $S$ for the growth rate, which would give a doubling time of 0.69 years. Growth is $S^2/2$, because the volatility drag takes half of the excess return at full Kelly.}
\end{solution}

\begin{solution}{iq:iv2:betting-and-kelly:23}
At leverage 2, growth is 10.88\% a year, about 87\% of the maximum of 12.5\%. Growth returns to zero at leverage 6.25, twice the optimum.

Growth is $g(f) = 0.08f - \tfrac12 \times 0.0256f^2 = 0.08f - 0.0128f^2$. At $f = 2$: $0.16 - 0.0512 = 0.1088$. As a multiple $c = 2/3.125 = 0.64$ of Kelly, the share is $c(2 - c) = 0.64 \times 1.36 \approx 0.87$. And $g = 0$ at $f = 0.08/0.0128 = 6.25$.

So a cap well below the optimum is close to costless. The limit takes 64\% of the risk and keeps 87\% of the growth.

\iqwrong{Treating 3.125 as the right leverage to run because it is the top. In this normal setting the top is flat, so going below it costs very little. Your risk of estimating $\mu$ too high, however, is large.}
\end{solution}

\begin{solution}{iq:iv2:betting-and-kelly:24}
The expected payoff is infinite. With a cap of $2^{20}$, it is \$21. A log investor with \$1,000 would pay about \$10.95.

Each $k$ contributes $2^{-k} \times 2^k = 1$, so the sum diverges: this is the St Petersburg paradox\index{St Petersburg paradox}. With the cap, the first 20 terms add to 20, and all longer runs together have probability $2^{-20}$ and pay $2^{20}$, which adds 1 more.

For log utility, solve $\sum_{k \geq 1}2^{-k}\ln(W - c + 2^k) = \ln W$ for $c$. At $W = 1{,}000$, this gives about 10.95. It rises slowly with wealth, to about 20.87 at \$1,000,000.

\iqwrong{Answering ``infinity, so any price''. Doubling the house's money adds only \$1 to the fair value, so the price is about the logarithm of what the house can actually pay.}
\end{solution}

\begin{solution}{iq:iv2:betting-and-kelly:25}
The standard error is about 0.5. Growth at full estimated Kelly is 0.375 a year if the true Sharpe ratio is 1, and $-0.375$ if it is 0.5. At half the estimated Kelly leverage, it is about 0.469 and about 0.094.

For roughly normal returns, the Sharpe ratio's standard error in annual terms is about $1/\sqrt{T}$, here $1/\sqrt 4$. Sized on the estimate $\hat S$, the volatility of the book is $\hat S$. If the true Sharpe ratio is $S$, growth per year is $\hat S S - \tfrac12\hat S^2$.

With $\hat S = 1.5$, growth is $1.5S - 1.125$. So $S = 1$ gives 0.375, and $S = 0.5$ gives $-0.375$. With half the size, the book's volatility is 0.75, and growth is $0.75S - 0.28125$. So $S = 1$ gives 0.469, and $S = 0.5$ gives 0.094.

An estimate of 1.5 with a standard error of 0.5 is quite compatible with a truth of 0.5. Full Kelly on that truth loses money, and half Kelly still compounds. Half Kelly also gives up little if the truth is 1.5: 0.84 against 1.125 a year. That is the argument for fractions of Kelly on estimated edges.

\iqwrong{Treating this as a matter of taste in risk. It holds even for a person who wants pure growth, because an estimated Sharpe ratio is noisy and usually too high.}
\iqfollow{Is 1.5 the right centre of the estimate?}{If you chose the best of many strategies, it is biased up. So a sizing rule should begin by shrinking it, and only then apply a fraction.}
\end{solution}

\begin{solution}{iq:iv2:betting-and-kelly:26}
Under full Kelly, $1/W_t$ is a martingale, and that bounds the chance of ever falling to $x$ by $x$. In continuous time, the probability at multiple $c$ is $x^{2/c - 1}$: $x$ at full Kelly and $x^3$ at half Kelly. So the chance of ever losing half is 50\% at full Kelly and 12.5\% at half.

For the discrete bound, with stake $f = 0.2$:
\[
E\left[\frac{1}{1 + fX}\right] = \frac{0.6}{1.2} + \frac{0.4}{0.8} = 1,
\]
so $M_t = 1/W_t$ is a positive martingale that starts at 1. For such a martingale, the chance of ever reaching a level $1/x$ is at most $x$, by the maximal inequality\index{Doob's maximal inequality} (or optional stopping). That is the bound. It is not an equality because $W$ jumps past the level: an exact calculation on the lattice gives about 45.5\% for $x = \tfrac12$, against 50\%. At the discrete half Kelly, it gives about 10.5\%.

In continuous time, $\ln W$ is Brownian motion with drift $m = (c - \tfrac12 c^2)S^2$ and volatility $v = cS$. The chance of ever going down by $L$ is $e^{-2mL/v^2}$, which equals $x^{2m/v^2}$ for $L = -\ln x$. Here $2m/v^2 = (2 - c)/c = 2/c - 1$.

\iqwrong{``Kelly never goes to zero, so it is safe.'' It never reaches zero, and yet there is a 50\% chance of a drawdown of a half and a 10\% chance of a drawdown of 90\%, from any start. Most people cannot live with that.}
\end{solution}

\begin{solution}{iq:iv2:betting-and-kelly:27}
The highest multiple is $c \approx 0.376$, and it keeps about 61\% of the maximum growth.

The chance of ever falling to a fraction $x$ is $x^{2/c - 1}$. Set this to 0.05 at $x = \tfrac12$: $2/c - 1 = \ln 0.05/\ln 0.5 \approx 4.32$, so $c = 2/5.32 \approx 0.376$. The share of maximum growth kept is $c(2 - c) = 0.376 \times 1.624 \approx 0.61$.

The derivation comes from the drift and volatility of log wealth. With drift $m = (c - \tfrac12 c^2)S^2$ and volatility $v = cS$, the probability that drifted Brownian motion ever goes $L$ below its start is $e^{-2mL/v^2}$. This does not depend on $S$ at all: only the multiple of Kelly matters, which is a nice thing to point out.

So a limit on drawdown is a statement in Kelly units, and investors who will not accept losing half are asking for about 0.38 Kelly.

\iqwrong{Trying to do this with a daily value-at-risk number. A limit on a drawdown over all time needs the hitting probability of the path, and not a one-day distribution.}
\iqfollow{And if the true Sharpe ratio is lower than you think?}{Then your real multiple $c$ is higher than you think, and the chance of a drawdown is higher. A $c$ of 0.376 on a Sharpe ratio estimated 50\% too high is really 0.564, and the chance of losing half rises to about 17\%.}
\end{solution}

\begin{solution}{iq:iv2:betting-and-kelly:28}
Stake about 19.6\%, which is $2E[p] - 1$. The posterior is $\mathrm{Beta}(61, 41)$, with mean about 0.598. Uncertainty about $p$ does not shrink the stake for a single bet, because expected log wealth depends on $p$ only through its mean.

Expected log wealth for stake $f$ is $E[p]\ln(1 + f) + (1 - E[p])\ln(1 - f)$, which is linear in $p$. So only the mean matters. The posterior standard deviation of about 0.048 plays no part.

Staking at the lower end, $p \approx 0.50$, means staking nothing. The argument against: under your stated model, you give up all of the expected growth, about 0.019, for no reason. Your model already includes the uncertainty.

The arguments for shrinking concern the model, not the arithmetic. The first is that a uniform prior is too generous: true edges are rare and small. Under a prior $\mathrm{Beta}(100, 100)$, the posterior mean is $160/300 \approx 0.533$, which gives a stake of about 6.7\%. The second is selection: this is a bettor you noticed because of a good run. The third is that the edge can change. So fix the prior, and then use the mean.

\iqwrong{Dividing by a variance of the estimate to ``penalise uncertainty''. This holds in some settings, but it has no basis in this one. Shrinkage should come from a better prior.}
\end{solution}

\begin{solution}{iq:iv2:betting-and-kelly:29}
The highest chance is 64.8\%, and you stake \$0.50 first. A Kelly 20\% each time gives 0\%, because three wins reach only \$1.728. Betting everything each time gives 21.6\%.

Work back from the end. With one bet left: at \$2 or more you are done, at \$1 to \$2 you can stake enough to reach \$2 with chance 0.6, and below \$1 the chance is 0. With two bets left: at \$1.50 or more, stake to reach \$2, and if that loses you are at \$1 or more with one bet left, so the chance is $0.6 + 0.4 \times 0.6 = 0.84$. From \$1 to \$1.50, the chance is 0.6. From \$0.50 to \$1, the chance is $0.6 \times 0.6 = 0.36$.

With three bets left at \$1, staking \$0.50 gives $0.6 \times 0.84 + 0.4 \times 0.36 = 0.648$. Other stakes do no better. This is dynamic programming\index{dynamic programming}, and it shows that the value function is a staircase, so the only stakes worth considering are those that land on a step.

\iqwrong{Using the Kelly stake because ``Kelly is optimal''. It maximises the logarithm of final wealth, an objective that is not the one you were given.}
\iqfollow{What if your objective were a smooth function of final wealth, not a hard target?}{The staircase goes, and the stake varies smoothly. A hard target is what produces these jumps, and also what makes people double up when they are behind.}
\end{solution}

\begin{solution}{iq:iv2:betting-and-kelly:30}
The rival has the higher expected wealth, $1.06^{1000}$ against $1.04^{1000}$. You finish ahead with probability about 95\%. Both hold: the rival's mean comes from paths that almost never happen.

Expected wealth per bet is $1 + 0.2f$, so more stake means a higher mean. To compare the paths, your ratio to the rival after $k$ wins in $n$ bets is $(1.2/1.3)^k(0.8/0.7)^{n - k}$. This is above 1 when $k/n < \ln(0.8/0.7)/\left(\ln(0.8/0.7) + \ln(1.3/1.2)\right) \approx 0.6252$. With 1,000 bets, that is at most 625 wins, and the binomial probability is about 0.95.

Typical log wealth is $1000g$: about 20.1 for you and about 14.7 for the rival. Log mean wealth is $1000\ln 1.04 \approx 39.2$ and $1000\ln 1.06 \approx 58.3$. Each mean is far above what either player typically sees. It is carried by paths in which the win share is high, and the rival wins on those.

\iqwrong{Saying that Kelly maximises expected wealth. Kelly gives up mean, and in return it beats any other fixed fraction with a probability that tends to 1 as the number of bets grows.}
\end{solution}

\begin{solution}{iq:iv2:betting-and-kelly:31}
The growth-optimal fraction is 50\%. Mean over variance gives about 89\%, and growth at that size is about $-0.039$ per bet, so you lose money over time on a bet with a positive edge.

Maximise $0.9\ln(1 + 0.25f) + 0.1\ln(1 - f)$. The derivative is zero when $0.225(1 - f) = 0.1(1 + 0.25f)$, so $0.125 = 0.25f$, and $f = 0.5$. There, growth is about 0.0367.

The mean is $0.9 \times 0.25 - 0.1 = 0.125$. The variance is $0.15625 - 0.125^2 \approx 0.1406$. The ratio is 0.889, and with the raw second moment in the denominator, it is 0.8. At 0.889, growth is $0.9\ln 1.2222 + 0.1\ln 0.1111 \approx -0.039$.

The approximation $\mu/\sigma^2$ comes from a second-order expansion, which is fine for small symmetric bets. It ignores skewness. Here a loss of the whole stake is a large move, and the logarithm punishes it harder than the expansion.

So sizing a strategy that sells options by its Sharpe ratio overstates its proper size, since the Sharpe ratio does not see what the crash does. Size from a stress loss: ask what happens to your wealth in the loss state, and take that logarithm.

\iqwrong{Taking a good Sharpe ratio as sufficient evidence. This trade's Sharpe ratio is 0.33 per bet, and it is a bad one to size by the formula.}
\end{solution}

\begin{solution}{iq:iv2:betting-and-kelly:32}
Before each card, stake the fraction $|R - B|/(R + B)$ on the colour that has more cards left. You end with $2^{52}/\binom{52}{26} \approx \$9.08$ on every order of the deck, with no risk at all. No strategy can have a higher expected final wealth.

Betting this way is the same as splitting your wealth in proportion to the chances: $R/(R + B)$ on red and $B/(R + B)$ on black. If a red card comes, wealth is multiplied by $2R/(R + B)$. Over the whole deck, the numerators give $26!\,26!$ whichever way the colours come, and the denominators give $52!$. So the multiplier is $2^{52}\,26!\,26!/52!$.

For the bound, note that an even-money stake moves wealth from one continuation to another without creating any. So summed over all $2^{52}$ colour sequences, final wealth totals $2^{52}$, whatever the strategy. Only $\binom{52}{26}$ sequences can occur, and wealth is not negative, so their average is at most $2^{52}/\binom{52}{26}$. The proportional strategy reaches it with certainty, and so by Jensen's inequality no strategy has a higher expected logarithm either.

\iqwrong{Waiting until the last card, where you know the colour for sure. That doubles your money once. The proportional bet earns on every card with an imbalance, and it has no variance.}
\iqwhy{It shows growth-optimal betting at its cleanest: when the bet's outcome leaves you with fewer possibilities, you are paid in information.}
\end{solution}

\begin{solution}{iq:iv2:betting-and-kelly:33}
The normal approximation needs about 1,075 bets at Kelly and about 476 at half Kelly. The exact binomial, taking the smallest number from which the chance stays at 95\% or more, gives 1,104 and 495.

To be ahead with 95\% confidence, you need $n g \geq 1.645\sqrt n\,s$, so $n \geq (1.645\,s/g)^2$, where $g$ is growth per bet and $s$ the standard deviation of log wealth. At $f = 0.10$: $g \approx 0.005008$ and $s \approx 0.0998$. At $f = 0.05$: $g \approx 0.003753$ and $s \approx 0.0498$.

Halving the stake cuts growth by a quarter but also cuts its standard deviation in half. So the ratio $g/s$ goes up by half, and the number of bets falls by more than half. For the exact count, you are ahead when wins exceed the fraction of bets at which log wealth is flat: $\ln(1/(1 - f))/\ln((1 + f)/(1 - f))$. You then take the binomial tail. The chance rises in a sawtooth with $n$, because wins are whole numbers, and that is why the exact count is a little higher.

\iqwrong{Assuming full Kelly reaches any goal fastest. It maximises long-run growth, but a goal in terms of confidence favours a smaller stake.}
\iqfollow{What stake is quickest to reach 95\%?}{As the stake goes to zero, $g/s$ tends to $(p - q)/(2\sqrt{pq})$ and the count falls to about 268, the same as the count for the bet's own sum to be positive. So smaller is better on this objective, and you pay in wealth reached.}
\end{solution}

\begin{solution}{iq:iv2:betting-and-kelly:34}
Final wealth is $2^n h!\,t!/(n + 1)! = 2^n/\left((n + 1)\binom{n}{h}\right)$, with $h$ heads and $t$ tails, whichever order they came in. With 7 heads in 10 tosses, you end with about \$0.776. Someone who knew $p = 0.7$ and played Kelly ends with about \$2.28. With 10 heads in 10, you end with about \$93.09, and the person who knew $p = 0.7$ ends with $1.4^{10}$, about \$28.93, so on that run you do better.

A stake $|2m - 1|$ on the likelier side is the same as splitting wealth into $m$ on heads and $1 - m$ on tails. So a head multiplies wealth by $2m$, and a tail by $2(1 - m)$. With a uniform prior, after $a$ heads and $b$ tails the posterior mean is $m = (a + 1)/(a + b + 2)$, known as Laplace's rule of succession\index{rule of succession}.

Now multiply along any path. The denominators are $2, 3, \ldots, n + 1$, which give $(n + 1)!$, in any order. Each head supplies the next of $1, 2, \ldots, h$, and each tail the next of $1, 2, \ldots, t$. So the product is $2^n h!\,t!/(n + 1)!$, and it depends on $h$ alone.

There is a cleaner view. Imagine splitting your dollar across all the possible $p$ in proportion to the prior, each piece playing Kelly with its $p$. The total after $n$ tosses is $\int_0^1 2^n p^h(1 - p)^t\,dp$, which is the same number. So Bayesian betting is a portfolio of Kelly bettors, and the best of them dominates in the long run.

Compare with the known $p = 0.7$: $2^{10} \times 0.7^7 \times 0.3^3 \approx 2.277$. That is 2.9 times what you made. In general, wealth falls short of the best fixed $p$ by a factor of at most $n + 1$, so the cost of learning grows only polynomially, while the wealth at a real edge grows exponentially.

\iqwrong{Thinking 7 heads out of 10 must make money. The early bets are made with little knowledge, and with so few tosses, learning costs more than the edge earns.}
\iqwhy{It gives a number to the cost of not knowing your edge. An edge you must learn is worth less, and it costs most at the beginning.}
\end{solution}

\begin{solution}{iq:iv2:betting-and-kelly:35}
The highest probability is $\Phi(1) \approx 84.1\%$. You achieve it by replicating a digital option that pays \$2 if the asset finishes above a level chosen so that the option costs exactly \$1. Kelly leverage gives about 42.3\%, and half Kelly about 26.2\%. If the asset has done badly late in the year, you hold very high leverage, because you have nothing to lose.

The market is complete, so any payoff costs its risk-neutral expectation. To finish with \$2 on an event $A$ and nothing elsewhere costs $2Q(A)$, so you need $Q(A) \leq \tfrac12$. You want to maximise $P(A)$ for that $Q(A)$. By the Neyman-Pearson lemma\index{Neyman-Pearson lemma}, take the states where $dP/dQ$ is largest, and these are the ones where the asset has risen most.

Let $B$ be Brownian motion under $P$. Under $Q$, $B^Q_t = B_t + St$ is Brownian motion. So $A = \{B^Q_1 \geq 0\}$ has $Q(A) = \tfrac12$, and $P(A) = P(B_1 + 1 \geq 0) = \Phi(1)$. In general, with wealth $w$, target $G$ and time $\tau$ left, the answer is $\Phi\left(\Phi^{-1}(w/G) + S\sqrt\tau\right)$. The asset's volatility does not appear, only the Sharpe ratio.

Wealth at time $t$ is the digital's value, $W_t = 2\Phi\left(B^Q_t/\sqrt{1 - t}\right)$, and the dollars of risk you run per unit of the asset's volatility are $\partial W/\partial B^Q = 2\varphi\left(B^Q_t/\sqrt{1 - t}\right)/\sqrt{1 - t}$. As a multiple of wealth, that is $\varphi(z)/\left(\Phi(z)\sqrt{1 - t}\right)$ with $z = B^Q_t/\sqrt{1 - t}$. It rises as $z$ falls and as the deadline nears, without limit. So if you are behind late in the year, you hold very high leverage.

For comparison, Kelly has $\ln W_1 = \tfrac12 + B_1$, and it needs $B_1 \geq \ln 2 - \tfrac12 \approx 0.193$, a chance of about 42.3\%. Half Kelly has $\ln W_1 = 0.375 + \tfrac12 B_1$, which needs $B_1 \geq 0.636$, a chance of about 26.2\%.

\iqwrong{Thinking a higher chance of success must be good. The losing 15.9\% of paths end at zero, a result neither Kelly nor half Kelly can give, and the best strategy uses unlimited leverage just before the deadline.}
\iqwhy{This is what a manager facing a hard target with a free option on the upside is pushed to do, and the reason for limits on leverage near the end of a period.}
\end{solution}
